% Folder 84, batch 06 — archivists' pages 101–120.
% Pass: Opus 5.5 (claude-opus-5-5), 2026-10-09 — first pass, unchecked against the pages by a human.
% Revised: Opus 5.5 (claude-opus-5-5), 2026-10-10 — re-read on the 700 dpi tiles; 34 \ill and 127 \uncertain resolved, 162 remain.
%
% Notation fixed for this batch: his cursive A (curly, sometimes close to a
% script letter) is plain A, A_1, A_2, A_3 throughout; the generator of A is
% \mathcal{U}, those of A_1, A_2 are U_1, U_2; the hook-shaped exponent he puts
% on A_1 \otimes A_2 and on its elements (pp. 101-104) is read as \natural;
% his contracted product is *; \Gamma for sections, \mathcal{O} for the base;
% \mathcal{X} for his script X (Spec A). The hand is fast throughout: the
% formulas carry, much of the connecting prose does not.
\documentclass[11pt,a4paper]{article}
\input{../preamble/grothendieck}

\folder{84}
\batch{6}
\pages{101}{120}
\foldertitle{[Formes quadratiques] : notes manuscrites (s.d.).}
\dating{[à partir de 1982-vers 1986]}
\watermark{Édition de démonstration}

\begin{document}

\page{101}
\note{la page reprend un calcul commencé avant le lot : la structure de $A_1 \otimes A_2$, les involutions et l'exposant $\natural$ (lecture douteuse d'un signe en crochet) sont introduits aux pages précédentes.}

\[
  \xi^{\natural} = \xi + \bar{\xi} \quad \text{pour } \xi \in A_1 \otimes A_2 ,
\]
\ill{} \ill{} \ill{} $\xi = 2$ i.e. \ill{}

\[
  x_1 \mapsto \bar{x}_1 = \text{\struck{$\xi_1(x_1)$}}\ \mathrm{Tr}_1(x_1) - x_1 ,
  \qquad
  x_2 \mapsto \bar{x}_2 = \mathrm{Tr}_2(x_2)\, e_2 - x_2
\]
\note{le second membre de $\bar{x}_2$ est écrit au-dessus d'une formule biffée illisible ; « $\mathrm{Tr}_1(x_1)$ » est ajouté au-dessus de « $\xi_1(x_1)$ » biffé.}
sur $A_1$, $A_2$, donc
\[
  \xi \mapsto \bar{\xi} , \qquad
  x_1 \otimes x_2 \mapsto \overline{x_1 \otimes x_2} \overset{\mathrm{déf}}{=} \bar{x}_1 \otimes \bar{x}_2
\]
sur $A_1 \otimes A_2$. On \uncertain{aura} donc, si $\xi \in A_1$, i.e.
$\xi = x_1 \otimes e_2$
\[
  \xi^{\natural} = (x_1 \otimes 1)^{\natural} = (x_1 + \bar{x}_1) \otimes e_2 = \mathrm{Tr}_1(x_1)\, e_1 \otimes e_2
\]
c'est OK.

\emph{NB} Si \struck{\ill{}} sur $X_0 = V(2 \cdot 1_X) \subset X$ (lieu de car.~2), \struck{\ill{}} \uncertain{tout} pt de $X_0$ l'une des deux structures $E_1$, $E_2$ est étale, donc $A \to A_1 \otimes A_2$ est univ. injectif, et \ill{} pour \ill{}
\[
  (A_1 \otimes A_2)^{\natural} = (A_1 \otimes A_2)^{\sigma_1 \otimes \sigma_2} .
\]
Ceci posé, on peut \add{donc} identifier l'application
\[
  (x_1, x_2) \mapsto x_1 * x_2 : A_1 \times A_2 \to A , \qquad
  x_1 \otimes x_2 \mapsto x_1 * x_2 : A_1 \otimes A_2 \to A
\]
\ill{} \ill{} des applications : à valeurs dans $(A_1 \otimes A_2)^{\natural}$. En tous cas \ill{} \ill{} \ill{} \uncertain{simplement} \struck{\ill{}} comme \emph{\uncertain{application}} trace pour $\sigma_1 \otimes \sigma_2$.

Quand $A_1$, $A_2$ sont des alg. \uncertain{de rev.} quadratiques, il s'agit bien de traces \uncertain{au sens habituel}. Mais cette trace, \ill{} \uncertain{cond.} de car.~$\neq 2$ \ill{},

\page{102}
\uncertain{Elle} \struck{\ill{}} sur $\mathcal{U}$, $U_1$, $U_2$ par
\[
  1 \mapsto 2 , \qquad U_1 \mapsto -b_1 , \qquad U_2 \mapsto -b_2 .
\]
\ill{} \ill{} $\Gamma\mathcal{O} \subset \Gamma A$, il \uncertain{suffit} de trouver l'image de $U_1 \otimes U_2$. C'est donc un élément de $A$ dont l'image dans $(A_1 \otimes A_2)^{\natural}$ est
\[
  U_1 \otimes U_2 + \underbrace{\bar{U}_1 \otimes \bar{U}_2}_{(-b_1 - U_1) \otimes (-b_2 - U_2)}
  = \underbrace{2 U_1 U_2 + b_1 U_2 + b_2 U_1}_{\text{image de } \mathcal{U}} + b_1 b_2 ,
\]
\struck{donc} il y a lieu de prendre
\[
  U_1 * U_2 = \mathcal{U} + b_1 b_2 .
\]

Mais il faut maintenant définir une application
\[
  N_{A_1 \otimes A_2 / A} : A_1 \otimes A_2 \to A ,
\]
application quadratique, dont la valeur en $z$ serait, essentiellement, $z \bar{z}$, mais « \uncertain{invariante} » \ill{} $A$. L'application bil. associée serait
\[
  (u, v) \mapsto u \bar{v} + v \bar{u} = \mathrm{Tr}_N(u \bar{v})
\]
déjà connue. Donc il faut se donner les valeurs de $N$, et de $\varphi_N$, sur les éléments de la base \uncertain{de} $A_1 \otimes A_2$. Or on \uncertain{devrait} \uncertain{avoir}, \ill{} $1$, $U_1$, $U_2$, $U_1 \otimes U_2$ \ill{}
\[
  N_A(x_1 \otimes x_2) = N_1(x_1)\, N(x_2) \in \Gamma\mathcal{O} \subset \Gamma A
\]
\[
  = (x_1 \otimes x_2)(\bar{x}_1 \otimes \bar{x}_2) = (x_1 \bar{x}_1) \otimes (x_2 \bar{x}_2) .
\]
En particulier
\[
  N(1) = 1 , \quad N(U_1) = N_1(U_1) = c_1 , \quad N(U_2) = N_2(U_2) = c_2 ,
\]
\[
  N(U_1 \otimes U_2) = c_1 c_2 .
\]

\page{103}
Peut \uncertain{se} définir comme \emph{à valeurs dans} $A$ (et non seulement comme à val. dans $(A_1 \otimes A_2)^{\natural}$, \uncertain{ce} qui ne serait pas si $A \to (A_1 \otimes A_2)^{\natural}$ \uncertain{est} \add{\ill{}} injectif, ni d'ailleurs surjectif).

On doit pouvoir définir \uncertain{de même} une norme, définie universellement par la condition que son image dans $(A_1 \otimes A_2)^{\natural}$ soit $\xi \bar{\xi}$.

Choisissons \uncertain{les} bases $(1, U_1)$ de $A_1$, $(1, U_2)$ de $A_2$, donc
\[
  A_1 = \mathcal{O}[U_1] / (U_1^2 + b_1 U_1 + c_1) , \qquad
  A_2 = \mathcal{O}[U_2] / (U_2^2 + b_2 U_2 + c_2) ,
\]
\uncertain{on trouve}
\[
  \begin{cases}
    A_1 \otimes A_2 = \mathcal{O}[U_1, U_2] / (U_1^2 + b_1 U_1 + c_1 ,\ U_2^2 + b_2 U_2 + c_2) \\
    A \simeq \mathcal{O}[\mathcal{U}] / (\mathcal{U}^2 + b\, \mathcal{U} + c)
  \end{cases}
\]
avec
\[
  \begin{cases}
    b = b_1 b_2 \\
    c = b_1^2 c_2 + b_2^2 c_1 - 4 c_1 c_2
  \end{cases}
\]
et $A \to A_1 \otimes A_2$ est donné par
\[
  \mathcal{U} \mapsto 2 U_1 U_2 + b_1 U_2 + b_2 U_1 .
\]

Explicitons
\[
  A_1 \otimes A_2 \xrightarrow{\ *\ } A , \qquad x_1 \otimes x_2 \mapsto x_1 * x_2 ,
\]
il suffit de trouver les images de $1$, $U_1 \otimes 1$, $1 \otimes U_2$, $U_1 \otimes U_2$.
\marginal{c'est donc ce qui tient lieu d'une trace de $A_1 \otimes A_2$ \uncertain{dans} $A$}
On aura bien sûr
\[
  x_1 * e_2 = \mathrm{Tr}_1(x_1) \cdot 1_A , \qquad
  e_1 * x_2 = \mathrm{Tr}_2(x_2) \cdot 1_A
\]
\marginal{ce qui définit $*$ sur $A_1 \oplus A_2 \subset A_1 \otimes A_2$}
\note{dans la première formule, « $e_2$ » est écrit au-dessus d'un symbole biffé ; dans la seconde, « $e_1$ » de même.}

\page{104}
\emph{Proposition.} Soient $A_1$, $A_2$ deux alg. quadratiques sur $\mathcal{O}$, et $A = A_1 \wedge A_2$. \struck{On a alors}.

1) \struck{\ill{}} \add{L'}homom. canonique
\[
  A_1 \otimes A_2 \xrightarrow{\ \mathrm{Tr}_{A_1 \otimes A_2 / A}\ } A
\]
(défini par la construction de $A$ comme ext. de $L_1^{\vee} \otimes L_2^{\vee}$ par $\mathcal{O}$, déduit de l'extension \add{$A_1 \otimes A_2$} de $L_1^{\vee} \otimes L_2^{\vee}$ par $A_1 \oplus_{\mathcal{O}} A_2 \subset A$, et de l'homom.
\[
  A_1 \oplus_{\mathcal{O}} A_2 \xrightarrow{\ \mathrm{Tr}\ } \mathcal{O}
\]
défini par $\mathrm{Tr}_{A_1}$ et $\mathrm{Tr}_{A_2}$ sur chaque summand) peut être considéré comme un « homom. trace ».
\note{sous la flèche, un « $*$ » est écrit sous l'indice de $\mathrm{Tr}$ ; dans la parenthèse, $A_1 \oplus A_2 \subset A$ porte en dessous la variante « $A_1 \oplus_{\mathcal{O}} A_2$ ».}

\marginal{NB On a \ill{} $\mathrm{Tr}\,(\mathrm{id}_{A_1} \otimes \sigma_2)(\xi)$ \ill{}, $\mathrm{Tr}\, \sigma_1 \otimes \mathrm{id}_{A_2}(\xi) = \mathrm{Tr}(\bar{\xi})$, $\mathrm{Tr}(\bar{\xi}) = \mathrm{Tr}(\xi)$ \uncertain{dès que} $\bar{\xi} = \sigma_1 \otimes \sigma_2(\xi)$}
\note{note marginale écrite en diagonale dans la marge gauche, en partie biffée.}

Son composé avec l'homom. can.
\[
  A \to A_1 \otimes A_2
\]
(qui d'ailleurs \uncertain{n'a} \uncertain{pas} à être univ. injectif, cf. plus haut) \struck{\ill{}} n'est autre que
\[
  \xi \mapsto \xi + \bar{\xi} , \qquad \text{i.e.} \qquad
  i\bigl(\mathrm{Tr}_{A_1 \otimes A_2 / A}\, \xi\bigr) = \xi + \bar{\xi}
\]
\[
  \bar{\xi} = \sigma_1 \otimes \sigma_2 (\xi)
\]
$\sigma_1$, $\sigma_2$ les involutions canoniques de $A_1$, $A_2$
\[
  x_1 \mapsto \mathrm{Tr}_{A_1}(x_1) - x_1 , \qquad x_2 \mapsto \mathrm{Tr}_{A_2}(x_2) - x_2 .
\]
\marginal{Donc en partic. $\mathrm{Tr}_{A_1 \otimes A_2 / A}(i(\eta)) = 2\eta$ pour \ill{}}

2) Il existe sur $A_1 \otimes A_2$ une application quadratique unique, à val. dans $A$, notée
\[
  N = N_{A_1 \otimes A_2 / A} ,
\]

\page{105}
ayant les propriétés suivantes
\begin{enumerate}
  \item[a)] $N(x_1 \otimes x_2) = N_{A_1}(x_1)\, N_{A_2}(x_2)$ \qquad $x_1 \in \Gamma A_1$, $x_2 \in \Gamma A_2$
  \item[b)] $\varphi_N(\xi, \eta) = \mathrm{Tr}_{A_1 \otimes A_2 / A}(\xi \bar{\eta})$
\end{enumerate}
\struck{On aura de} i.e.
\[
  N(\xi + \eta) = N(\xi) + N(\eta) + \mathrm{Tr}(\xi \bar{\eta})
\]
On aura alors, pour tt $\xi \in A_1 \otimes A_2$
\[
  \begin{cases}
    i(N(\xi)) = \xi \bar{\xi} \\
    N(\sigma_1 \otimes \mathrm{id}_{A_2}(\xi)) = N(\mathrm{id}_{A_1} \otimes \sigma_2(\xi)) = \overline{N(\xi)} \\
    N(\bar{\xi}) = N(\xi)
  \end{cases}
\]
\marginal{Notamment $N_{A_1 \otimes A_2 / A}(i(\eta)) = \eta^2$}
\note{dans la deuxième ligne, l'indice de $\mathrm{id}$ est écrit « $A_2$ » dans les deux termes ; on lit $\mathrm{id}_{A_1}$ dans le second, comme l'exige le sens.}

De plus, on a
\[
  N_{A/\mathcal{O}}\bigl(N_{A_1 \otimes A_2 / A}(\xi)\bigr) = N_{A_1 \otimes A_2 / \mathcal{O}}(\xi)
\]
\uncertain{ainsi que} la formule plus bénigne
\[
  \mathrm{Tr}_{A/\mathcal{O}}\bigl(\mathrm{Tr}_{A_1 \otimes A_2 / A}(\xi)\bigr) = \mathrm{Tr}_{A_1 \otimes A_2 / \mathcal{O}}(\xi)
\]

Prouvons p.ex. la première \uncertain{et} \uncertain{la} deuxième. \ill{} Posons
\[
  \sigma_1 = \sigma_{A_1} \otimes \mathrm{id}_{A_2} , \qquad
  \sigma_2 = \mathrm{id}_{A_1} \otimes \sigma_{A_2} , \qquad
  \sigma_3 = \sigma_1 \sigma_2
\]
\struck{deuxième} \struck{donc} $\xi = \xi_0 \in A$, \struck{on doit vérifier} $\xi_i = \sigma_i(\xi)$, on constate qu'on ait
\[
  N_{A_1 \otimes A_2 / \mathcal{O}}(\xi) = \xi_0\, \xi_1\, \xi_2\, \xi_3 .
\]
\note{« $\xi_0 \in A$ » : le $A$ est celui qu'il emploie pour le produit ; on attendrait $A_1 \otimes A_2$. Sous $\xi_0 \xi_1 \xi_2 \xi_3$, des accolades groupent $\xi_1 \xi_2$ et $\xi_0 \cdots \xi_3$.}

\page{106}
Or le deuxième membre est
\[
  \xi_0 \bar{\xi}_0\, \xi_1 \bar{\xi}_1 , \qquad \bar{\xi}_0 = \xi_3 , \quad \bar{\xi}_1 = \xi_2 ,
\]
or
\[
  \xi_0 \bar{\xi}_0 \overset{!}{=} N_{A_1 \otimes A_2 / A}(\xi_0)
\]
\[
  \xi_1 \bar{\xi}_1 = N_{A_1 \otimes A_2 / A}(\xi_1) = \overline{N_{A_1 \otimes A_2 / A}(\xi_0)} ,
  \qquad \xi_1 = \sigma_1(\xi_0)
\]
donc \struck{\ill{}} \uncertain{si} $N = N_{A_1 \otimes A_2 / A}(\xi_0)$, on trouve
\[
  N \bar{N} = N_{A/\mathcal{O}}(N) \qquad \text{OK.}
\]

Quand on axiomatise la situation \struck{de $A$, $A_1$, \ill{}} de $\mathrm{Tr}_{A_1 \otimes A_2 / A}$ et $N_{A_1 \otimes A_2 / A}$, on aboutit à la situation d'une \ill{} sur $k$, (\add{d'}une $k$-algèbre $A$ qui « se comporte comme si c'était une algèbre quadratique sur $k$ » (i.e. loc. libre de rang~2). Mais on ne suppose ni que $A$ soit loc. libre \add{ni de rang 2} \struck{\ill{}}/$k$, ni que
\[
  \alpha : k \to A , \qquad \lambda \mapsto \lambda \cdot 1_A
\]
soit injective. Au lieu de cela, on suppose données deux applications
\[
  t_{A/k} = t : A \to k \quad (\text{trace relative}) , \qquad
  n_{A/k} = n : A \to k \quad (\text{norme relative}) ,
\]
satisfaisant les conditions suivantes

(a) $t$ est $k$-linéaire, $n$ est quadratique, et $t$ se déduit de $n$, comme
\[
  t(x) = \varphi_n(x, 1) \qquad (\text{la forme bil. ass. à } n)
\]
donc \ill{}
\[
  n(x + \lambda 1_A) = n(x) + \lambda^2 n(1) + \lambda t(x) , \qquad n(1) = 1 \ \text{cf. ci-dessous b)}
\]
\note{devant la formule encadrée, un mot biffé.}

\page{107}
\struck{D'autre part, la \ill{}} \emph{NB} On \uncertain{verra} plus \uncertain{loin} comme la forme bil. qui s'explicite en termes de \struck{la} forme $n$.

(b) \quad $n(1) = 1$, \quad $n(xy) = n(x)\, n(y)$ \quad i.e. $n$ est un homom. de monoïdes unitaires $A^{\times} \to k^{\times}$

(c) \quad $t(1) = 2$
\note{la lettre (c) est écrite sur un (b) ; à côté de la condition encadrée, une ligne biffée commençant par « D'où » et un mot biffé (« Posons »).}

Posons $\tau(f) = \alpha(t(f))$, et considérons l'application
\[
  \sigma : f \mapsto \bar{f} = \tau(f) - f
\]
\struck{Ceci \uncertain{suppose} que cette application \ill{} \ill{}}
On a
\[
  \tau(\bar{f}) = \tau(f) \underbrace{\tau(1)}_{2} - \tau(f) = \tau(f)
\]
d'où
\[
  (\bar{f})^{-} = \tau(\bar{f}) - \bar{f} = f
\]
i.e. \struck{$\sigma$ est} $\sigma^2 = \mathrm{id}_A$.
\marginal{Notons que par définition $\tau(f) \overset{\mathrm{déf}}{=} \alpha(t(f)) = f + \bar{f}$, et aussi $\bar{1} = 1$}

Ceci posé, on veut que l'on ait

1°) $\sigma$ est un automorphisme d'algèbre, i.e.
\[
  \sigma(xy) = \sigma(x)\, \sigma(y) \qquad x, y \in A
\]
ce qui s'explicite comme
\[
  2xy = \underbrace{\tau(xy)}_{\varphi_n(x, \bar{y})} - \tau(x)\tau(y) + x\tau(y) + y\tau(x)
\]
\note{« $-\tau(x)\tau(y)$ » est ajouté au-dessus de la ligne ; sous $\tau(xy)$, « $= \varphi_n(x,\bar{y})$ $\to$ voir ci-dessous ».}
\marginal{comb. lin. à coeff. dans $k$ de $1$, $x$, $y$}

\emph{NB} Le facteur 2 ici (qui n'est gênant que \struck{en car.~2} si 2 \ill{} \ill{} dans $k$), \uncertain{on voit} donc que la structure \ill{} d'algèbre de $A$ est connue quand on connaît le $k$-module, avec l'élément $1_A$ ($A$ \uncertain{en tant que} $k$-module, avec $\lambda \mapsto \lambda \cdot 1_A$) et des \add{$\alpha : k \to A$} applic. $t : A \to k$, $n : A \to k$.

\marginal{\emph{NB} Quand $\alpha$ est injectif, \ill{} \ill{} \ill{} ⑤ \ill{} $n(1) = 1$ \ill{} $\alpha\, n(x) = x\bar{x}$ \ill{} \uncertain{$\alpha\, n(xy) = n(x)\, n(y)$}}
\note{note marginale en diagonale, en grande partie illisible.}

2°)
\[
  \alpha(n(x)) = x \bar{x}
\]
implique $\alpha(n(x)) = \alpha(n(\bar{x}))$, mais on \ill{} \ill{} \uncertain{qui} \struck{\ill{}} \uncertain{soit}
\[
  n(\bar{x}) = n(x)
\]
On a
\[
  n(\bar{x}) = n(\lambda 1_A - x) \overset{(a)}{=} n(x) + \lambda^2 - \lambda t(x) \qquad \text{où } \lambda = t(x)
\]
donc \uncertain{on a} \uncertain{bien} $n(\bar{x}) = n(x)$.
\note{dans la formule, des termes biffés au-dessus de « $\lambda^2$ ».}

\page{108}
La formule ci-dessus pour \struck{\ill{}}
\[
  \alpha(n(x)) = x \bar{x}
\]
donne, si on fait $x = y + z$
\[
  \underbrace{\alpha(n(y+z))}_{n(y) + n(z) + \varphi_n(y,z)}
  = \text{\struck{$\ldots$}}\ (y+z)\overline{(y+z)}
\]
\[
  = \underbrace{y\bar{y}}_{\alpha(n(y))} + \underbrace{z\bar{z}}_{\alpha(n(z))}
  + \underbrace{y\bar{z} + z\bar{y}}_{\tau(y\bar{z})} ,
  \qquad z\bar{y} = \overline{y\bar{z}}
\]
donc
\[
  \alpha(\varphi_n(y, z)) = \tau(y \bar{z}) \qquad \text{i.e. } \alpha(t(y\bar{z}))
\]
donc si $\alpha$ est injectif, on en déduit
\[
  \varphi_n(y, z) = t(y \bar{z})
\]
i.e. la formule d'additivité pour $n$
\[
  n(y + z) = n(y) + n(z) + t(y\bar{z}) , \qquad t(y\bar{z}) = t(\bar{y} z)
  \quad \text{car } t(x) = t(\bar{x})
\]

Ceci posé, on veut

3°) la formule précédente est vraie (sans supposer $\alpha$ injectif…)

4°) Une formule facultative, quand $k$ est fini loc. libre \struck{de rang} \struck{\ill{}} $n$ sur un anneau $k_0$, \uncertain{et} $A$ loc. libre de rang $2n$ \struck{et \ill{} \ill{}} sur $k_0$ (N.B. \struck{\ill{}} $k_0 \to A$ injectif mais $k \to A$ injectif)\note{tel quel sur la page ; le sens demande « mais non $k \to A$ injectif ».}. \uncertain{On} \uncertain{veut} \uncertain{supposer} \uncertain{que}
\[
  \begin{cases}
    \mathrm{Tr}_{k/k_0}(t(x)) = \mathrm{Tr}_{A/k_0}(x) \\
    N_{k/k_0}(n(x)) = N_{A/k_0}(x)
  \end{cases}
\]
\marginal{identité entre applic. linéaires $A \to k_0$ ; identité de $(2n)$-formes}

\page{109}
Récapitulons le formulaire pour

\begin{tikzcd}
  k \arrow[r, "\alpha"] & A \arrow[l, bend left=40, "t"] \arrow[l, bend left=70, "n"]
\end{tikzcd}

($k$-alg.)
\marginal{$\alpha$ homom. d'anneaux comm. $\alpha(\lambda) = \lambda \cdot 1_A$. On pose $\tau(x) = \alpha\, t(x)$, $\nu(x) = \alpha\, n(x)$}

\note{ici et pp.~111, 113, les numéros (1) à (6) sont cerclés.}

(1) $t$ \struck{k}-$k$-linéaire, $t(1_A) = 2$

(2) $x \mapsto \bar{x} \overset{\mathrm{déf}}{=} \tau(x) - x : A \to A$ \quad (automorphisme involutif du $k$-module $A$, avec $\bar{1}_A = 1_A$) est un automorphisme d'alg., i.e.
\[
  \overline{xy} = \bar{x}\,\bar{y} \qquad \text{i.e.} \qquad
  2xy = \bigl(\tau(xy) - \tau(x)\tau(y)\bigr) + \tau(y)\,x + \tau(x)\,y
\]

(3)
\[
  \begin{cases}
    n(\lambda x) = \lambda^2 n(x) \\
    n(x + y) = n(x) + n(y) + t(x\bar{y})
  \end{cases}
  \qquad \lambda \in k , \ x \in A
\]
i.e. $n$ est quadratique, la forme bil. ass. est
\[
  \varphi_n(x, y) = t(x \bar{y})
\]

(4) $n(1) = 1$ \quad (\emph{NB} compte tenu de la formule d'additivité en 3°, cela implique $t(1) = 2$)

(5) $\alpha(n(x)) = x \bar{x}$ \quad $x \in A$ \quad (qui implique 3° 4° si $\alpha$ injectif)

(6) (facultatif), \add{et} si $k$ loc. libre de rang $n$ sur $k_0$, et $A$ loc. libre de rang $2n$ sur $k_0$)
\[
  \begin{cases}
    \mathrm{Tr}_{k/k_0}(t(x)) = \mathrm{Tr}_{A/k_0}(x) \\
    N_{k/k_0}(n(x)) = N_{A/k_0}(x)
  \end{cases}
  \qquad x \in A
\]

La donnée de \add{$A$ sur $k$, et} $\alpha$, $t$, $n$ satisfaisant 1° à 5° sera appelée algèbre \emph{quasi-quadratique} sur $k$. On va maintenant pouvoir étendre dans la situation \add{quasi-}\emph{biquadratique}.

\page{110}
Données
\begin{enumerate}
  \item[a)] un anneau de base $k$
  \item[b)] $k$-algèbres $A_1$, $A_2$, $A_3$, $A$
  \item[c)] Sur chacune des $A_i$, structure quasi-quadratique sur $k$, i.e. on a
\end{enumerate}

\begin{tikzcd}
  A_i \arrow[r, bend left=20, "t_i"] \arrow[r, bend right=20, "n_i"'] & k
\end{tikzcd}

\struck{\uncertain{Comme pour}} satisfaisant les conditions 1° à 5° ci-dessus. Les notat. $x \mapsto \bar{x}$ l'involution canonique de $A_i$,
\[
  \bar{x} = t_i(x) - x
\]
\struck{$\bar{x} = t_i$} (\add{par abus de} notation) (\emph{NB} on identifie $\lambda \in k$ avec ses images dans les $A_i$, sans pour autant supposer $k \to A_i$ injectif…)
\begin{enumerate}
  \item[d)] Des homom. de $k$-algèbres
\end{enumerate}

\begin{tikzcd}
  A_1 \arrow[dr, "\alpha_1"'] & A_2 \arrow[d, "\alpha_2"] & A_3 \arrow[dl, "\alpha_3"] \\
  & A &
\end{tikzcd}

(non néc. inj.), et des structures d'alg. quasi-quadratiques de $A$ sur les $A_i$

\begin{tikzcd}[column sep=large, row sep=large]
  & A \arrow[dl, bend right=15, "T_1"'] \arrow[dl, bend left=15, "N_1"] \arrow[d, bend right=15, "T_2"'] \arrow[d, bend left=15, "N_2"] \arrow[dr, bend right=15, "T_3"'] \arrow[dr, bend left=15, "N_3"] & \\
  A_1 & A_2 & A_3
\end{tikzcd}

\note{une première version du diagramme, à droite, est biffée.}
données satisfaisant donc au formulaire ci-dessus de 1° à 5°. De plus, \emph{compatibilités}.

\struck{On pose}
\[
  \text{\struck{$\sigma_i(x) = T_i(x) - x$}}
\]
\struck{$\sigma_i$ est un auto. involutif de $A$ ; et on veut que}
\[
  \text{\struck{$\sigma_1 \sigma_2 = \sigma_3 , \quad \sigma_2 \sigma_3 = \sigma_1 , \quad \sigma_3 \sigma_1 = \sigma_2$}}
\]
\struck{i.e. que les $\sigma_i$ sont des involutions qui \emph{commutent} entre elles, et $\sigma_1 \sigma_2 \sigma_3 = \mathrm{id}$}
\note{tout ce dernier passage est barré de traits obliques.}

\page{111}
Considérons un des 6 diagrammes \uncertain{carrés}.

\begin{tikzcd}
  A_1 \arrow[r, no head, "\alpha_1"] \arrow[d, no head] & A \arrow[d, no head, "\alpha_2"] \\
  k \arrow[r, no head] & A_2
\end{tikzcd}

On veut que la situation soit, du point de vue des \add{\ill{}} \uncertain{traces}, \ill{}, comme si $A_1$ était quadratique sur $k$, et que $A$ en soit déduit par \uncertain{changement} de base $k \to A_2$.
\marginal{Mais ici on \uncertain{suppose} $A/A_2$ ($\alpha_2$) quasi-quadratique \ill{} \ill{} $A_1$, $A_2$ \ill{}}
\note{note marginale en diagonale, presque entièrement illisible ; les traits du carré sont tracés sans pointes.}

(1)
\[
  T_2\bigl(\alpha_1(x_1)\, \text{\struck{$\alpha_2(x_2)$}}\bigr) = t_1(x_1)\, \text{\struck{$\alpha_2(x_2)$}}
\]
\marginal{ce qui détermine $T_2$ sur $k[A_1, A_2] \subset A$ en termes de $t_1$}

i.e. on veut un diagramme \uncertain{commutatif}.

\begin{tikzcd}
  A_1 \otimes_k A_2 \arrow[r] \arrow[d, "t_1 \otimes \mathrm{id}_{A_2}"'] & A \arrow[d, "T_1"] \\
  A_2 \arrow[r, no head] & A_2
\end{tikzcd}

\note{la flèche verticale de droite est lue $T_1$ ; d'après (1) on attendrait $T_2$. En bas, les deux $A_2$ sont reliés par un trait double.}
\marginal{NB. $A_1 \otimes A_2 \to A$ n'est pas supposé injectif, mais que $t_1 \otimes \mathrm{id}_{A_2}$ s'annule \uncertain{au moins} sur le noyau}

Ceci implique que
\[
  \sigma_2(\alpha_1(x)) = \alpha_1(\bar{x}) ,
\]
i.e. $\sigma_2$ (\uncertain{ainsi} \uncertain{que} $\sigma_3$) « \uncertain{induit} » sur $A_1$ l'involution donnée $x_1 \mapsto \bar{x}_1$ de $A_1$.

On veut une condition similaire pour les normes

2°
\[
  N_2(\alpha_1(x_1)) = n_1(x_1)
\]
\note{dans les deux membres, un facteur biffé illisible suit $\alpha_1(x_1)$ et $n_1(x_1)$.}

\page{112}
Je dis que cette formule implique que le diagramme similaire au précédent

\begin{tikzcd}
  A_1 \otimes_k A_2 \arrow[r] \arrow[d, "{n_1 \otimes_k \mathrm{id}_{A_2} = N_{12}}"'] & A \arrow[d, "N_2"] \\
  A_2 \arrow[r, no head] & A_2
\end{tikzcd}

\note{en bas, les deux $A_2$ sont reliés par un trait double fléché.}
\marginal{déduite de la forme $n_1$ par chgt de base $k \to A_2$}
commute. C'est en tous cas vrai pour les éléments \struck{de la forme} de $A_1 \otimes_k A_2$ de la forme $x_1 \otimes x_2$, car les deux membres sont \uncertain{alors}
\[
  t_1(x_1) \cdot x_2^2
\]
\note{le premier facteur se lit $t_1(x_1)$ ; le sens demande $n_1(x_1)$. Entre les deux facteurs, des lettres biffées.}
par quadraticité de $N_{1,2}$, $N_2$. Il reste à voir qu'on a commutativité pour les \struck{formes} formes bil. associées, or pour $N_{12}$ c'est la forme déduite de
\[
  \text{\struck{$(x,y) \mapsto$}}\ \varphi_{n_1}(x, y) = t_1(x \bar{y}) \quad \text{sur } A_1 \times A_1 ,
\]
par chgt. de base $k \to A_1$, or on a, si $x' = \alpha_1(x)$, $y' = \alpha_1(y)$
\[
  \begin{cases}
    \alpha_1(\bar{y}) = \sigma_2(\alpha_1(y)) \\
    t_1(x \bar{y}) = T_2(\alpha_1(x\bar{y})) \quad \text{dans } A_2
  \end{cases}
\]
\note{« $k \to A_1$ » est écrit tel ; on attendrait $k \to A_2$.}
donc
\[
  \varphi_{n_1}(x, y) = T_2\bigl(x' \sigma_2(y')\bigr) = \varphi_{N_2}(x', y') \qquad \text{OK.}
\]

\page{113}
\emph{NB} La condition sur $N_2$ implique celle sur $T_2$, en passant aux formes bil. ass., i.e. en \uncertain{l'appliquant} pour $x$ et pour $1 + x$ ($x \in A_1$).

(3) \uncertain{Ensuite} on veut la condition
\[
  \sigma_1 \sigma_2 \sigma_3 = 1
\]
qui implique que les involutions $\sigma_i$ commutent entre elles. Car en prenant 2 d'entre elles on \uncertain{a une} involution
\[
  \begin{cases}
    \sigma_1 \sigma_2 = \sigma_3 \\
    \sigma_2 \sigma_3 = \sigma_1 \\
    \sigma_3 \sigma_1 = \sigma_2
  \end{cases}
\]
(car pour deux invol. $\sigma$ et $\tau$,
\[
  \sigma\tau = \tau\sigma \iff (\sigma\tau)^2 = 1 \quad \text{i.e. } \sigma\tau\sigma\tau = 1
\]
la condition $\sigma_1 \sigma_2 \sigma_3 = 1$ \ill{})
\marginal{Cette condition est automatiquement vérifiée sur la sous-algèbre de $A$ engendrée par les $A_i$, car elle est vraie sur les $A_i$ (\uncertain{lesquels} \ill{} \ill{} $\sigma_i$) \uncertain{et} \uncertain{commutent}}
\note{note marginale en diagonale, à gauche de (3).}

En utilisant 1°), \struck{\ill{}} s'explicite par la formule
\[
  2x = T_1(x) + T_2(x) + T_3(x) - t_2 T_2(x) ,
\]
\uncertain{qui} n'est pas symétrique, \add{mais} elle est équivalente à
\[
  \sigma_2 \sigma_3 \sigma_1 = 1 \quad \text{ou} \quad \sigma_3 \sigma_1 \sigma_2 = 1 ,
\]
donc
\[
  2x = T_1(x) + T_2(x) + T_3(x) - t_3 T_3(x)
\]
\[
  2x = \cdots - t_1 T_1(x) ,
\]
elle implique (*) la condition
\[
  t_1 T_1(x) = t_2 T_2(x) = t_3 T_3(x)
\]
(qui ressemble à 6° ci-dessus). Si on désigne par $T_{A/k}(x)$ cette valeur commune aux trois membres, on trouve donc la condition
\[
  2x = T_1(x) + T_2(x) + T_3(x) - T_{A/k}(x) .
\]
\marginal{(*) donc \ill{} si $k \to A$ injectif, mais \ill{} \ill{} vérifiée \uncertain{en tous cas}\ldots}

\page{114}
\note{tout le haut de la page, jusqu'à la seconde formule encadrée, est barré de traits obliques.}
\struck{Sous réserve que les $\alpha_i : A_i \to A$ soient bien injectifs, cela montre que chacun des $T_i$ est déterminé par les deux autres, \uncertain{via}}
\[
  \text{\struck{$\alpha_1 T_1(x) = 2x + T_{A/k}(x)$}}
\]
\[
  \text{\struck{${} - \alpha_2 T_2(x) - \alpha_3 T_3(x)$}}
\]
\struck{On s'attend à trouver des formules similaires, multiplicatives}
\[
  \text{\struck{$N_{A/k}(x)\, x^2 = N_1(x)\, N_2(x)\, N_3(x)$}}
\]
\struck{et}
\[
  \text{\struck{$n_1(N_1(x)) = n_2(N_2(x)) = n_3(N_3(x))$}}
\]

Mais \struck{toutes} les \struck{conditions} \add{relatives sur la norme} (à prouver ci-dessous), sont conséquences de celles sur les traces, i.e. de $\sigma_1 \sigma_2 \sigma_3 = 1$ sur les $\sigma_i$. En effet, posons
\[
  x_0 = x , \quad x_1 = T_1(x) , \quad x_2 = T_2(x) , \quad x_3 = T_3(x)
\]
\note{les lettres sont lues $T_i$ ; le sens et la suite demandent $x_i = \sigma_i(x)$.}
On a alors
\[
  \begin{aligned}
    \sigma_1 &: x_0 \leftrightarrow x_1 , \quad x_2 \leftrightarrow x_3 \\
    \sigma_2 &: x_0 \leftrightarrow x_2 , \quad x_1 \leftrightarrow x_3 \\
    \sigma_3 &: x_0 \leftrightarrow x_3 , \quad x_1 \leftrightarrow x_2
  \end{aligned}
\]
d'où \struck{$n_1(N_1(x)) = N_1(x)\overline{N_1(x)} = \ldots$, $N_1(x_1) = \ldots$}
\[
  n_1(N_1(x)) \overset{!}{=} \underbrace{N_1(x)}_{x\, \sigma_1(x)} \overline{N_1(x)}
  = (x_0 x_1)\, \underbrace{\sigma_2(x_0 x_1)}_{x_2 x_3} = x_0 x_1 x_2 x_3
\]
et de même pour $n_2(N_2(x))$, $n_3(N_3(x))$. Donc on a l'égalité des $n_i(N_i(x))$ si $k \to A$ est injectif. Et on a
\[
  N_1(x)\, N_2(x)\, N_3(x) = (x_0 x_1)(x_0 x_2)(x_0 x_3)
\]
\[
  = x_0^2 \underbrace{(x_0 x_1 x_2 x_3)}_{N_{A/k}(x_0)} = x^2 N_{A/k}(x) ,
\]
OK.

\page{115}
Considérons l'algèbre produit tensoriel \uncertain{triple}
\[
  A_1 \otimes A_2 \otimes A_3 ,
\]
via les $\sigma_{A_i}$, le groupe $\mathbb{F}_2^3$ opère dessus. Il n'y a pas en général d'opération de ce groupe sur $A$, compatible avec
\[
  \textstyle\bigotimes A_i \to A ,
\]
car il faudrait qu'il y ait p.ex. un auto. $\tau_1$ de $A$ qui induise \struck{l'} l'identité sur $A_2$, $A_3$, et $\sigma_{A_1}$ sur $A_1$. Mais dans le cas le plus simple, $A$ est engendré par $A_2$, $A_3$, donc $\tau_1$ devrait être l'identité, donc qu'il induise $\sigma_{A_1}$ sur $A_1$ ne l'est pas.

Par contre, soit
\[
  \Gamma \subset \mathbb{F}_2^3 , \qquad
  \Gamma = \mathrm{Ker}\bigl(\mathbb{F}_2^3 \to \mathbb{F}_2\bigr) , \quad (\varepsilon_i) \mapsto \textstyle\sum \varepsilon_i
\]
formé des éléments $(\varepsilon_1, \varepsilon_2, \varepsilon_3)$ tels que le nombre de composantes $\neq 0$ (i.e. $= 1$) soit $0$ ou $2$ i.e. soit pair. Les \struck{\ill{}} trois élts non nuls de ce groupe, savoir
\[
  \text{\struck{$(0,0,1)$}}, \quad (0,1,1) , \quad (1,0,1) , \quad (1,1,0)
\]
qui correspondent aux automorphismes
\[
  \begin{aligned}
    \mathrm{id}_{A_1} \otimes \sigma_{A_2} \otimes \sigma_{A_3} &\overset{\mathrm{déf}}{=} \tau_1 \\
    \sigma_{A_1} \otimes \mathrm{id}_{A_2} \otimes \sigma_{A_3} &\overset{\mathrm{déf}}{=} \tau_2 \\
    \sigma_{A_1} \otimes \sigma_{A_2} \otimes \mathrm{id}_{A_3} &\overset{\mathrm{déf}}{=} \tau_3
  \end{aligned}
\]

\page{116}
et \struck{ici} on a une \struck{\ill{}} op. de $\Gamma$ sur $A$, telle que
\[
  \textstyle\bigotimes A_i \to A
\]
soit compatible aux deux opérations de $\Gamma$, savoir celle qui fait opérer $\tau_i$ par $\sigma_{A/A_i}$.

\struck{Considérons} Supposons maintenant que les $A_i$ correspondent \struck{à} des rev. quadratiques $X_i$ de $X$, d'où \uncertain{alors} un rev. quadratique
\[
  X_1 \wedge X_2 \wedge X_3 , \qquad \text{d'algèbre } B = A_1 \wedge A_2 \wedge A_3 ,
\]
et \uncertain{on a}
\[
  X_1 \times X_2 \times X_3 \to X_1 \wedge X_2 \wedge X_3 , \qquad
  (x_1, x_2, x_3) \mapsto x_1 * x_2 * x_3 ,
\]
i.e.
\[
  B \to A_1 \otimes A_2 \otimes A_3 ,
\]
et de façon plus précise, on a
\[
  B \to (A_1 \otimes A_2 \otimes A_3)^{\Gamma} \qquad \text{invariant sous } \Gamma
\]

Et même, si on exclut les pts $x \in X = \mathrm{Spec}(k)$ de \emph{car.~2} tels que deux au moins des $A_i$ soient non étales en ce point, dans l'homom. précédent est un \emph{iso}.

On \uncertain{va} \uncertain{regarder} le composé
\[
  B \to A_1 \otimes A_2 \otimes A_3 \to A , \quad \text{\uncertain{i.e.}}
\]
\note{la phrase continue p.~117.}

\page{117}
i.e., en termes de $\mathcal{X} = \mathrm{Spec}\, A$, le composé
\[
  \mathcal{X} \to X_1 \times_X X_2 \times_X X_3 \to X_1 \wedge X_2 \wedge X_3 ,
\]
j'ai envie de prouver que c'est une \uncertain{flèche} \emph{constante}, i.e. qu'il se factorise par une \emph{section} de $X_1 \wedge X_2 \wedge X_3$, ce qui signifie aussi que l'homom. composé précédent

\begin{tikzcd}
  B \arrow[rr, "f"] \arrow[dr] & & A \\
  & (\bigotimes A_i)^{\Gamma} \arrow[ur] &
\end{tikzcd}

est à valeurs dans $k \cdot 1_A$. (Du moins, on y \uncertain{aboutira} quand on suppose $k \to A$ injectif — c'est ce qu'on va supposer maintenant, pour simplifier.)

Mais on a évidemment
\[
  f(B) \subset A^{\Gamma} ,
\]
et la condition est vérifiée si on suppose
\[
  A^{\Gamma} = k \cdot 1_A .
\]
(C'est là une condition pas toujours vérifiée, \uncertain{même} dans des cas \uncertain{raisonnables}, à cause de la car.~2. Mais si 2 est inversible dans $k$, c'est toujours vrai sous les conditions précédentes).

\page{118}
En effet, si $x \in A$, $\sigma_i(x) = x$ signifie \struck{\ill{}}
\[
  T_i(x) - x = x \quad \text{i.e.} \quad 2x = T_i(x) ,
\]
et si c'est vrai pour $i = 1, 2, 3$ on trouve
\[
  6x = \underbrace{T_1(x) + T_2(x) + T_3(x)}_{= 2x + T_{A/k}(x) \ \text{cf. plus haut}}
\]
d'où
\[
  4x = T_{A/k}(x)
\]
ce qui implique, si 2 inv.,
\[
  x = \tfrac{1}{4} T_{A/k}(x) \in k \qquad \text{OK.}
\]

\struck{\ill{} \ill{} \ill{} \ill{} \ill{} \ill{} $A^{\Gamma} = k \cdot 1_A$}
\note{deux lignes barrées de croix.}

\emph{Proposition.} Supposons les $A_i$ des alg. quadratiques \emph{étales}, que $k \to A$ soit injectif, que les $\alpha_i(A_i) \subset A$ engendrent l'algèbre $A$ i.e. $\bigotimes A_i \to A$ épi. (donc
\[
  \mathcal{X} = \mathrm{Spec}\, A \to X_1 \times_k X_2 \times_k X_3
\]
est une immersion) (\uncertain{enfin} $A^{\Gamma} = k \cdot 1_A$, \ill{} \ill{}). Alors \struck{la situation \ill{}} \add{\ill{}}
\struck{se déduit de la situation \ill{}}
\[
  \mathcal{X} \to X_1 \times X_2 , \quad \mathcal{X} \to X_2 \times X_3 , \quad \mathcal{X} \to X_3 \times X_1
\]
sont des iso., $\mathcal{X}$ est \add{fini} \emph{étale} sur $X = \mathrm{Spec}\, k$.

Et la situation se déduit par \emph{twist} de celle \struck{\ill{}} $I_1$, $I_2$, $I_3$, ens. à deux élts munis d'une trivialisation $\varepsilon$ de $I_1 \wedge I_2 \wedge I_3$, et de
\[
  J \subset I_1 \times I_2 \times I_3 , \qquad
  J = \{ (a_1, a_2, a_3) \in I_1 \times I_2 \times I_3 \mid a_1 \wedge a_2 \wedge a_3 = \varepsilon \}
\]

\page{119}
Inversement, \add{dans} une telle situation tordue, $A$ est une algèbre étale quadratique sur chacun des $A_i$, donc on a une situation biquadratique \add{étale}, et de plus les hypothèses ci-dessus (en particulier $A^{\Gamma} = k \cdot 1_A$) sont satisfaites.

\struck{Il suffit} Seule la première partie, directe, demande un peu de soin. OPS $X_1$, $X_2$, $X_3$ \ill{} $I_{1X}$, $I_{2X}$, $I_{3X}$, et \add{la section constante} \struck{la section} \add{rev. trivialisé de $X_1 \wedge X_2 \wedge X_3 = (I_1 \wedge I_2 \wedge I_3)_X$} \uncertain{constante}. \struck{\ill{}} \uncertain{On a} alors
\[
  \mathcal{X} \subset J_X
\]
et tout résulte du lemme suivant

\emph{Lemme.} Soit $\Gamma$ un groupe, opérant de façon transitive sur un ens. $J$, et considérons un sous-schéma \add{fermé} $\mathcal{X}$ de $X_J$ (où $X = \mathrm{Spec}\, k$) tel que
\begin{enumerate}
  \item[a)] $\mathcal{X}$ est stable sous $\Gamma$, et
  \item[b)] l'homomorphisme $k \to \Gamma(X_J, \mathcal{O}_{X_J})$ $(\simeq k^J)$ est injectif.
\end{enumerate}
\note{en b), on attendrait l'anneau de $\mathcal{X}$ plutôt que celui de $X_J$ ; transcrit tel quel.}
Alors $\mathcal{X} = X_J$, donc $\mathcal{X}$ est étale sur $X$ (c'est fini étale si $J$ est fini).

En effet, soit $j \in J$, considérons le composant $X_j$ de $X_J$, $X_j \xrightarrow{\ \sim\ } X$, et $\mathcal{X}_j = X_j \cap \mathcal{X}$, il correspond à un sous-schéma fermé de $X$ donc à un idéal $\mathfrak{J}_j$ de $k$. À cause de l'\uncertain{action} de

\page{120}
$\mathcal{X}$ sous $\Gamma$, et de la transitivité de $\Gamma$ sur $J$, on voit que cet idéal $\mathfrak{J}_j$ ne dépend pas de $j$. D'autre part, l'intersection des $\mathfrak{J}_j$ est évidemment le noyau de
\[
  k \to \Gamma(X_J, \mathcal{O}_{X_J}) \simeq k^J .
\]
Comme celui-ci est nul, on a $\mathfrak{J}_j = 0$ \uncertain{pour tout} $j$, \uncertain{donc}
\[
  \mathcal{X}_j = X_j \qquad \text{donc} \qquad \mathcal{X} = X_J \qquad \text{OK.}
\]

Ce résultat vaut sans restrictions de caractéristiques, mais il concerne la situation standard étale, \emph{très} restrictive. J'ai envie de vérifier que dès qu'on a trois rev. quadratiques $X_1$, $X_2$, $X_3$ de $X$, et une section $\varepsilon$ de $X_1 \wedge X_2 \wedge X_3$, sans condition étale, on en déduit sur
\[
  \mathcal{X} = \pi^{-1}(\varepsilon) \qquad
  \bigl(\pi : X_1 \times_X X_2 \times_X X_3 \to X_1 \wedge X_2 \wedge X_3\bigr)
\]
une structure biquadratique par rapport aux $X_i$. Ceci doit généraliser la structure biquadratique connue sur un $X_1 \times X_2$, par les applications

\begin{tikzcd}[row sep=small]
  & X_1 \\
  X_1 \times X_2 \arrow[ur] \arrow[r] \arrow[dr] & X_2 \\
  & X_3 = X_1 \wedge X_2
\end{tikzcd}

comme on \uncertain{va le voir} de suite.
\note{l'argument se poursuit au-delà du lot.}

\end{document}
